% !TeX program = xelatex
\documentclass[UTF8,12pt,a4paper]{ctexart}

% ---------- 常用宏包 ----------
\usepackage[a4paper,margin=2.5cm]{geometry}
\usepackage{amsmath,amssymb,amsthm,mathtools,bm}
\usepackage{enumitem}
\usepackage{fancyhdr}
\usepackage{lastpage}
\usepackage{xcolor}
\usepackage{hyperref}

% ---------- 请修改这里 ----------
\newcommand{\coursename}{工程硕士数学}
\newcommand{\homeworknumber}{HW1}
\newcommand{\studentname}{陈天乐}
\newcommand{\studentid}{[student ID removed]}
\newcommand{\classinfo}{深微硕6班}
\newcommand{\submissiondate}{\today}

% ---------- 页面设置 ----------
\setlength{\parindent}{2em}
\setlength{\parskip}{0.35em}
\linespread{1.25}

\pagestyle{fancy}
\fancyhf{}
\lhead{\coursename}
\chead{\homeworknumber}
\rhead{\studentname}
\cfoot{第 \thepage\ 页，共 \pageref{LastPage} 页}
\renewcommand{\headrulewidth}{0.4pt}

\hypersetup{
  colorlinks=true,
  linkcolor=blue,
  urlcolor=blue,
  pdftitle={\coursename—\homeworknumber},
  pdfauthor={\studentname}
}

% ---------- 题目与解答环境 ----------
\newcounter{problem}
\newenvironment{problem}[1][]
  {\refstepcounter{problem}\par\medskip
   \noindent\textbf{\large 题目 \theproblem\if\relax\detokenize{#1}\relax\else\quad #1\fi}
   \par\smallskip\noindent}
  {\par\medskip}

\newenvironment{solution}
  {\noindent\textbf{解：}\quad}
  {\hfill$\square$\par\medskip}

\newenvironment{proofcn}
  {\noindent\textbf{证明：}\quad}
  {\hfill$\blacksquare$\par\medskip}

% 常用命令
\newcommand{\R}{\mathbb{R}}
\newcommand{\norm}[1]{\left\lVert #1\right\rVert}
\newcommand{\abs}[1]{\left\lvert #1\right\rvert}
\DeclareMathOperator{\rank}{rank}
\DeclareMathOperator{\diag}{diag}

\begin{document}

% ---------- 标题 ----------
\begin{center}
  {\LARGE\bfseries \coursename}\par
  \vspace{0.4em}
  {\Large \homeworknumber}\par
  \vspace{1em}
  \begin{tabular}{rl@{\qquad}rl}
    姓名：& \studentname & 学号：& \studentid \\
    班级：& \classinfo   & 日期：& \submissiondate
  \end{tabular}
\end{center}

\vspace{0.5em}
\hrule
\vspace{1em}

% ---------- 正文示例 ----------
\begin{problem}[]
已知 \(f(x)=\sin x,\,x\in[0,2\pi]\)。试求 \(C[0,2\pi]\) 中函数的范数 \(\|f\|_\infty\) 和 \(\|f\|_2\)
\end{problem}

\begin{solution}
\[
  \|f\|_\infty=\max\limits_{x\in[0,2\pi]}|f(x)|=1
\]

\[
\|f\|_2=\sqrt{\int_0^{2\pi}f^2(x) dx} = \sqrt{\int_0^{2\pi}\sin^2(x) dx }= \sqrt{\int_0^{2\pi}\frac{1-\cos(2x)}{2} dx}=\sqrt{\pi}
\]
\end{solution}

\begin{problem}
求向量 \(\|x\|_\infty,\|x\|_1,\|x\|_2\)\\
(1) \(\boldsymbol{x}=(2,1,-3,4)^T\)\\
(2) \(\boldsymbol{x}=(\sin k,\cos k,2^k)^T,\ k\in\mathbb{N}\)
\end{problem}

(1)
\begin{solution}
\[
  \|x\|_\infty =\max\{|2|,|1|,|-3|,|4|\}=4
\]
\[
  \|x\|_1 =|2|+|1|+|-3|+|4|=2+1+3+4=10
\]
\[
  \|x\|_2 =\sqrt{2^2+1^2+(-3)^2+4^2}=\sqrt{4+1+9+16}=\sqrt{30}
\]
\end{solution}

(2)
\begin{solution}
\[
\|x\|_\infty=\max\left\{|\sin k|,|\cos k|,|2^k|\right\}=2^k 
\]

\[
\|x\|_1=|\sin k|+|\cos k|+2^k
\]

\[
\|x\|_2=\sqrt{\sin^2k+\cos^2k+(2^k)^2}=\sqrt{1+4^k}
\]

\end{solution}

\newpage
\begin{problem}[]
证明：对一切 \(\boldsymbol{x}\in\mathbb{R}^n\)，有\\
(1) \(\|\boldsymbol{x}\|_\infty \le \|\boldsymbol{x}\|_1 \le n\|\boldsymbol{x}\|_\infty\)；\\
(2) \(\|\boldsymbol{x}\|_\infty \le \|\boldsymbol{x}\|_2 \le \sqrt{n}\|\boldsymbol{x}\|_\infty\)。

\end{problem}

\begin{proofcn}

(1)

不妨设\(i=i_0\)时
\[
\max\limits_i |x_i|=|x_{i_0}|=\|x\|_\infty
\]

那么
\[
\|x\|_1=\sum_{i=1}^n |x_i| = |x_{i_0}|+\sum_{i=1,i\neq i_0}^n |x_i| \ge |x_{i_0}| = \|x\|_\infty
\]

并且易得
\[
\|x\|_1=\sum_{i=1}^n |x_i| = |x_{i_0}|+\sum_{i=1,i\neq i_0}^n |x_i| \le n\times |x_{i_0}| = n\|x\|_\infty
\]


(2)


不妨设\(i=i_0\)时
\[
\max\limits_i |x_i|=|x_{i_0}|=\|x\|_\infty
\]

那么
\[
\|x\|_2=\sqrt{(\sum_{i=1}^n |x_i|^2)} = \sqrt{(|x_{i_0}|^2+\sum_{i=1,i\neq i_0}^n |x_i|^2)}\ge\sqrt{|x_{i_0}|^2} = |x_{i_0}|=\|x\|_\infty
\]

并且易得
\[
\|x\|_2=\sqrt{(\sum_{i=1}^n |x_i|^2)} = \sqrt{(|x_{i_0}|^2+\sum_{i=1,i\neq i_0}^n |x_i|^2)}\le\sqrt{n\times|x_{i_0}|^2} = \sqrt{n}|x_{i_0}|=\sqrt{n}\|x\|_\infty
\]


\end{proofcn}

\newpage

\begin{problem}[]
求下列矩阵 \(\boldsymbol{A}\) 的范数 \(\|\boldsymbol{A}\|_1\)，\(\|\boldsymbol{A}\|_2\)，\(\|\boldsymbol{A}\|_\infty\) 以及 \(\rho(\boldsymbol{A})\)\\
(1)\(\boldsymbol A=\begin{bmatrix}1 & -2\\-3 & 4\end{bmatrix}\)\\
(2)\(\boldsymbol A=\begin{bmatrix}2 & -1 & 0\\-1 & 2 & -1\\0 & -1 & 2\end{bmatrix}\)

\end{problem}

\begin{solution}

(1)

\[
\|A\|_1=\max\limits_{1\le j\le 2} \sum_{i=1}^{2}|A_{ij}|=6
\]

\[
det(\lambda I-A) = |\lambda I-A|=\begin{vmatrix}\lambda-1 & 2\\3 & \lambda-4\end{vmatrix}=(\lambda-1)(\lambda-4)-6=\lambda^2-5\lambda-2=0
\]


解得
\[
\rho(A)=|\lambda_{max}|=\frac{5+\sqrt{33}}{2}
\]

\[
A^TA=
\begin{bmatrix}
1 & -3\\
-2 & 4
\end{bmatrix}
\begin{bmatrix}
1 & -2\\
-3 & 4
\end{bmatrix}
=
\begin{bmatrix}
10 & -14\\
-14 & 20
\end{bmatrix}
\]

\[
det(\lambda I-A^TA)=|\lambda I-A^TA|=\begin{vmatrix}\lambda-10 & 14\\14 & \lambda-20\end{vmatrix}=(\lambda-10)(\lambda-20)-196=\lambda^2-30\lambda+4=0
\]

解得
\[
\|A\|_2 = \sqrt{\lambda_{max}}=\sqrt{\frac{30+2\sqrt{221}}{2}}=\sqrt{15+\sqrt{221}}
\]

\[
\|A\|_\infty=\max\limits_{1\le i\le 2} \sum_{j=1}^{2}|A_{ij}|=7
\]


(2)

\[
\|A\|_1=\max\limits_{1\le j\le 2} \sum_{i=1}^{2}|A_{ij}|=4
\]

\[
det(\lambda I-A) = |\lambda I-A|=\begin{vmatrix}\lambda-2 & 1 & 0\\1 & \lambda-2 & 1 \\ 0 & 1 & \lambda-2\end{vmatrix}=(\lambda-1)(\lambda-4)-6=\lambda^2-5\lambda-2=0
\]

令 \(t=2-\lambda\)，则
\[
\det(A-\lambda I)
=
\begin{vmatrix}
t&-1&0\\
-1&t&-1\\
0&-1&t
\end{vmatrix}
=t^3-2t
\]

因此
\[
t(t^2-2)=0,
\]

即
\[
t=0,\quad t=\sqrt2,\quad t=-\sqrt2.
\]

因此
\[
\rho(A)=|\lambda_{max}|=2+\sqrt{2}
\]

观察到\(A\)是实对称方阵，那么其存在正交对角化形式
\[
A=Q\Lambda Q^T\qquad \Lambda=\operatorname{diag}(\lambda_1,\lambda_2,\lambda_3)
\]

因此
\[
A^TA=Q\Lambda^2Q^T
\]

故
\[
\|A\|_2=\sqrt{\lambda_{max}(A^TA)}=|\lambda_{max}(A)|=2+\sqrt{2}
\]
\[
\|A\|_\infty=\max\limits_{1\le i\le 3} \sum_{j=1}^{3}|A_{ij}|=4
\]

\end{solution}




\begin{problem}[]
已知 \(\boldsymbol{A},\boldsymbol{B}\in\mathbb{R}^{n\times n}\)，且 \(\det\boldsymbol{A}\neq0,\det\boldsymbol{B}\neq0\)。证明：对任意一种从属的矩阵范数\\
(1) \(\|\boldsymbol{A}^{-1}\| \ge \dfrac{1}{\|\boldsymbol{A}\|}\);\\
(2) \(\|\boldsymbol{A}^{-1}-\boldsymbol{B}^{-1}\| \le \|\boldsymbol{A}^{-1}\|\,\|\boldsymbol{B}^{-1}\|\,\|\boldsymbol{A}-\boldsymbol{B}\|\)。

\end{problem}

\begin{proofcn}

(1)

欲证明该命题，由矩阵范数正定性，只需证
\[
\|A^{-1}\| \|A\| \ge 1 
\]

对\(\forall\boldsymbol{x}\in\mathbb R^n\)，由相容性
\[
\|A^{-1}A \boldsymbol{x}\|=\|\boldsymbol{x}\|\le \|A^{-1}\| \|A\| \|\boldsymbol{x}\|
\]

两边同时消掉\(\|\boldsymbol{x}\|\)，得证

(2)

对\(\forall\boldsymbol{x}\in\mathbb R^n\)，观察到
\[
\|(A^{-1}-B^{-1}) \boldsymbol{x}\|=\|(B^{-1}-A^{-1}) \boldsymbol{x}\|=\|B^{-1}(A-B)A^{-1}\boldsymbol{x}\|
\]

由相容性
\[
\|B^{-1}(A-B)A^{-1})\boldsymbol{x}\| \le \|B^{-1}\|\|A-B\|\|A^{-1}\|\|\boldsymbol{x}\|
\]

\[
\therefore \frac{\|B^{-1}(A-B)A^{-1})\boldsymbol{x}\|}{\|\boldsymbol{x}\|} = \frac{\|(A^{-1}-B^{-1}) \boldsymbol{x}\|}{\|\boldsymbol{x}\|} \le \|A^{-1}\|\|B^{-1}\|\|A-B\|
\]

由\(\boldsymbol{x}\)的任意性可知
\[
\max\limits_{i} \frac{\|(A^{-1}-B^{-1}) \boldsymbol{x}\|}{\|\boldsymbol{x}\|}  =\|A^{-1}-B^{-1}\| \le \|A^{-1}\|\|B^{-1}\|\|A-B\|
\]

\end{proofcn}

\begin{problem}[]
已知 \(\boldsymbol{P}\in\mathbb{R}^{n\times n}\)，\(\det\boldsymbol{P}\neq0\)。\(\|\boldsymbol{x}\|\) 为 \(\mathbb{R}^n\) 上一种向量范数，从属于它的矩阵范数为 \(\|\boldsymbol{A}\|\)。\\
(1) 定义 \(\|\boldsymbol{x}\|_P = \|\boldsymbol{P}\boldsymbol{x}\|\)，试证 \(\|\cdot\|_P\) 是 \(\mathbb{R}^n\) 上的一种向量范数；\\
(2) 证明从属于向量范数 \(\|\boldsymbol{x}\|_P\) 的矩阵范数 \(\|\boldsymbol{A}\|_P = \|\boldsymbol{P}\boldsymbol{A}\boldsymbol{P}^{-1}\|\)。
\end{problem}

\begin{proofcn}

  (1)

  \(\because \|\cdot\|\)是向量范数，由向量范数的正定性,对\(\forall \boldsymbol{x}\in \mathbb{R}^{n}\)
  \[
    \|\boldsymbol{x}\|_P=\|P\boldsymbol{x}\|\ge 0
  \]

  满足正定性。同时，若\(\|\boldsymbol{x}\|_P=\|P\boldsymbol{x}\|=0\)且\(\boldsymbol{P}\)非奇异
  \[
  \therefore\|P\boldsymbol{x}\|=0\,\mbox{当且仅当}\,\boldsymbol{x}=0 
  \]

  对\(\forall \alpha \in \mathbb{C}\)
\[
\|\alpha \boldsymbol{x}\|_P=\|\alpha P\boldsymbol{x}\|=|\alpha|\| P\boldsymbol{x}\|=|\alpha|\|\boldsymbol{x}\|_P
\]

满足齐次性

  对\(\forall \boldsymbol{x},\boldsymbol{y} \in \mathbb{R}^n\)
\[
\| \boldsymbol{x}+\boldsymbol{y} \|_P=\|P(\boldsymbol{x}+\boldsymbol{y})\|=\|P\boldsymbol{x}+P\boldsymbol{y}\|\le \|P\boldsymbol{x}\|+\|P\boldsymbol{y}\|=\|\boldsymbol{x}\|_P+\|\boldsymbol{y}\|_P
\]


满足三角不等式。综上，得证\(\|\boldsymbol{x}\|_P\)是一种向量范数


  (2)
\[
\|A\|_P=\max \frac{\|A\boldsymbol{x}\|_P}{\|\boldsymbol{x}\|_P}=\max \frac{\|PA\boldsymbol{x}\|}{\|P\boldsymbol{x}\|}
\]
  
令\(\boldsymbol{y}=P\boldsymbol{x}\)。因为\(P\)是可逆矩阵，那么\(\boldsymbol{x}=P^{-1}\boldsymbol{y}\)。上式可写成


\[
\|A\|_P=\max \frac{\|PAP^{-1}\boldsymbol{y}\|}{\|\boldsymbol{y}\|}=\|PAP^{-1}\|
\]
  

\end{proofcn}


\end{document}
